\section{Generador de carga, pruebas y despliegue}

Este capítulo cierra el manual con la cara operativa del sistema: cómo se le
inyecta carga para medirlo, cómo se prueba su lógica con un arnés en Java puro y
cómo se despliega y se opera el clúster sobre Google Cloud el día de la
evaluación.

\subsection{Generador de carga}

El generador de carga vive en \codigo{loadtest/LoadDriver.java}, una clase
\emph{stand-alone} que es el único ``cliente'' del proyecto: no depende de nada
del servidor más allá del contrato REST publicado. Genera tráfico contra un solo
nodo (el líder) durante una duración fija mezclando un \textbf{80\% de lecturas
de saldo y un 20\% de transferencias}, proporción gobernada por la constante
\codigo{READ\_RATIO = 0.80}. Por defecto corre \codigo{DEFAULT\_SECONDS = 60}
segundos con \codigo{DEFAULT\_CLIENTS = 8} hilos concurrentes, pero los
argumentos de \codigo{main} permiten fijar segundos, hilos y el rango de cuentas
\codigo{[minId, maxId]}; subiendo los segundos puede correr de forma
prácticamente indefinida, como exige la prueba del profesor que ``nunca
termina''.

Antes de medir, \codigo{bootstrapToken()} registra un usuario desechable
(\codigo{loaddriver}) y hace \codigo{login} para obtener el JWT que comparten
todos los hilos; tanto un \codigo{201} (creado) como un \codigo{409} (ya
registrado) son aceptables porque solo se necesita poder iniciar sesión. Un
detalle de rendimiento importante es que \emph{un} único \codigo{HttpClient}
(\codigo{sharedClient}) se reparte entre todos los \emph{workers}: un cliente por
hilo abriría cientos de hilos selectores propios y el generador se ahogaría en su
propio \emph{context-switch} antes de saturar al banco.

Cada hilo es un \codigo{Worker} que, hasta una fecha límite común en
\codigo{System.nanoTime()}, sortea la operación según \codigo{READ\_RATIO} y la
despacha. Solo se cuentan los éxitos reales: \codigo{doRead()} incrementa
\codigo{reads} únicamente si el servidor devolvió \codigo{HTTP 200} con un saldo
parseable, y \codigo{doTransfer()} incrementa \codigo{transfers} solo si la
transferencia fue aceptada (también \codigo{200}).

\begin{lstlisting}[language=Java,caption={Bucle de carga por hilo (Worker.run, recortado)}]
public void run() {
    HttpClient http = sharedClient;
    try {
        while (System.nanoTime() < deadlineNanos) {
            if (ThreadLocalRandom.current().nextDouble() < READ_RATIO) {
                doRead(http);
            } else {
                doTransfer(http);
            }
        }
    } finally {
        done.countDown();
    }
}
\end{lstlisting}

La parte didácticamente más valiosa es la verificación de la invariante de
conservación del dinero. Antes de arrancar la carga,
\codigo{captureTotalCents()} lee \codigo{GET /api/stats} una sola vez y guarda
\codigo{totalBalance} (en centavos): leer el total que el servidor suma en
memoria es O(1) en el cliente, en lugar de sumar cientos de miles de saldos por
HTTP. Al terminar, \codigo{verifyAndReport()} vuelve a leer el total y exige que
sea \textbf{idéntico} al inicial: ninguna transferencia crea ni destruye dinero,
así que la suma global debe permanecer constante. Si no coincide,
\codigo{reportCulprit()} recorre el rango buscando la cuenta con saldo negativo
para señalar al culpable; si coincide, imprime \codigo{CONSISTENT} junto con las
lecturas, las transferencias y el \codigo{score = transferencias*4 + lecturas}
(una transferencia vale cuatro veces una lectura, según el PDF del concurso).

\begin{lstlisting}[language=Java,caption={Verificacion de la invariante de conservacion (verifyAndReport, recortado)}]
void verifyAndReport(long initialTotalCents, String scenario) {
    long finalTotalCents = captureTotalCents();
    if (finalTotalCents != initialTotalCents) {
        System.out.println("INCONSISTENT: expected " + initialTotalCents
                + " cents, found " + finalTotalCents + " cents");
        reportCulprit(initialTotalCents, finalTotalCents);
        return;
    }
    long score = transfers.get() * 4 + reads.get();
    System.out.println("CONSISTENT");
    // ...
}
\end{lstlisting}

El cuerpo de la transferencia respeta la forma de cable obligatoria:
\codigo{transferBody()} serializa \codigo{sourceAccountId} y
\codigo{targetAccountId} como cadenas y \codigo{amount} como decimal; los montos
salen de \codigo{randomAmount()} en el rango \codigo{[0.01, 100.00]}, valores
bajos a propósito para que las transferencias rara vez agoten una cuenta y el
tráfico no degenere en puros rechazos por saldo insuficiente.

\subsection{Harness de pruebas}

El proyecto no usa JUnit: trae su propio arnés de pruebas en Java puro, formado
por dos piezas. \codigo{tests/Assert.java} es un ayudante mínimo de aserciones
(\codigo{isTrue}, \codigo{equals}, \codigo{notNull}, \codigo{throwsException})
donde cada método lanza un \codigo{AssertionError} con un mensaje descriptivo
cuando la expectativa falla. \codigo{tests/RunTests.java} es el punto de entrada:
instancia cada \emph{suite}, ejecuta sus casos y, capturando el
\codigo{AssertionError}, marca cada caso como \codigo{PASS} o \codigo{FAIL} sin
detener a los demás. Un caso es un \codigo{record Case(String name, Runnable
body)}, de modo que la falla se expresa simplemente lanzando una excepción desde
el cuerpo. Al final imprime un resumen \codigo{total/passed/failed} y termina con
estatus distinto de cero si algo falló, de manera que sirve para frenar la
compilación.

\begin{lstlisting}[language=Java,caption={Recoleccion y ejecucion de suites (RunTests.main, recortado)}]
List<Case> cases = new ArrayList<>();
cases.addAll(new LedgerTest().cases());
cases.addAll(new MoneyTest().cases());
cases.addAll(new RequestParserTest().cases());
cases.addAll(new WireCodecTest().cases());
cases.addAll(new AuthTest().cases());
cases.addAll(new ReplayTest().cases());
cases.addAll(new CheckpointTest().cases());
cases.addAll(new JwtCrossNodeTest().cases());
cases.addAll(new CatchupTest().cases());
// ... corre cada caso, cuenta PASS/FAIL y hace System.exit(1) si hubo fallos
\end{lstlisting}

Las \emph{suites} están organizadas por subsistema. \codigo{tests/MoneyTest.java}
verifica que el dinero se redondee y convierta correctamente entre cadena decimal
y centavos (con \codigo{HALF\_UP}) y que el viaje de ida y vuelta sea exacto.
\codigo{tests/LedgerTest.java} comprueba que \codigo{Ledger.move} transfiera
centavos exactos entre dos cuentas y mantenga el total constante bajo
concurrencia. \codigo{tests/RequestParserTest.java} valida que el parser arme una
petición tanto en una sola entrega como partida entre varias.
\codigo{tests/WireCodecTest.java} cubre la codificación de una transferencia a una
línea JSON y su \emph{round-trip}. \codigo{tests/AuthTest.java} prueba la
autorización del \codigo{Authenticator}: el esquema \codigo{Bearer} es
insensible a mayúsculas (RFC 7235) y tolera espacios, mientras que un encabezado
ausente, un esquema equivocado o un token vacío o falso se rechazan.

El bloque de mayor valor para este sistema distribuido son las pruebas de
recuperación y de réplica. \codigo{tests/ReplayTest.java} fija el contrato de
\codigo{Bank.applyReplicated()}: un seguidor reproduce en orden de secuencia y debe
\emph{forzar} la aplicación de cada transferencia (ya fue autorizada por el líder)
en vez de re-verificar fondos, porque re-verificar podría rechazar y
\emph{descartar} una transferencia autorizada mientras los contadores avanzan,
divergiendo de forma permanente; aplicar incondicionalmente es independiente del
orden y por tanto convergente. La \emph{suite} alimenta una secuencia en orden
adverso, comprueba la convergencia y que un reenvío de una secuencia ya vista sea
idempotente. \codigo{tests/CheckpointTest.java} cubre las
primitivas de \emph{checkpoint} que permiten a una réplica reanudar desde su
secuencia exacta tras un reinicio total: \codigo{Bank.snapshotInto} (captura
atómica de marca de agua y saldos), \codigo{Bank.seedLastSeq} (siembra la marca
desde un \emph{checkpoint} restaurado) y la garantía de que un \emph{snapshot}
restaurado más el delta del líder converge exactamente a los saldos por cuenta
del líder. \codigo{tests/CatchupTest.java} fija la reanudación por secuencia exacta:
que \codigo{log.since(X)} sea justamente la cola \codigo{X+1..} y que un banco
sembrado en la marca \codigo{X} descarte los reenvíos \codigo{<=X} y aplique desde
\codigo{X+1}. \codigo{tests/JwtCrossNodeTest.java} fija la propiedad entre nodos del
JWT: un token emitido con el secreto compartido valida en otro nodo y uno firmado
con un secreto distinto se rechaza.

Existen además dos verificaciones de integración que \textbf{no} corre
\codigo{RunTests} porque abren un socket real:
\codigo{tests/ReplConsistency.java} cablea un \codigo{Bank} líder con uno réplica
sobre \codigo{ReplicaFeed} y \codigo{ReplicaSync}, lo martilla con varios hilos
escritores haciendo miles de transferencias concurrentes y luego exige que el
saldo \textbf{por cuenta} líder $=$ réplica para cada cuenta (la suma global se
conserva aun si una transferencia se pierde o se duplica, así que solo la
comparación por cuenta atrapa un \emph{bug} de replicación). \codigo{tests/ReplInversion.java}
es la regresión de la convergencia por cuenta bajo carga concurrente con saldos
deliberadamente ajustados y montos de hasta el saldo completo -- la ruta que
originalmente expuso el \emph{bug} de inversión de secuencia --, que
\codigo{ReplConsistency} no ejercita porque sus cuentas arrancan demasiado holgadas. Cada una imprime \codigo{REPL\_CONSISTENCY\_OK} o
\codigo{REPL\_INVERSION\_OK} y sale con \codigo{0} en éxito; se ejecutan a mano
con el \emph{fat jar} en el \emph{classpath}.

\subsection{Despliegue en GCP}

La infraestructura se aprovisiona con \codigo{deploy/crear-infra.sh}, que crea en
Google Cloud (zona \codigo{us-central1-a}) una IP estática para el líder, las
reglas de \emph{firewall} (puerto \codigo{8080} abierto a clientes y
\emph{dashboard}, \codigo{9090} solo entre nodos con la etiqueta
\codigo{tesoreria}) y exactamente \textbf{2x \codigo{e2-standard-2} + 1x
\codigo{e2-standard-4}}, respetando la restricción de usar solo la serie E2. Los
tres nodos arrancan el \emph{mismo} \codigo{deploy/startup-script.sh} y difieren
únicamente por metadatos por nodo: \codigo{tes-node-id}, \codigo{tes-leader-host},
\codigo{tes-peers} y un \codigo{tes-jwt-secret} compartido. La convención clave es
que un \codigo{tes-leader-host} \emph{vacío} marca al nodo como líder (que recibe
la IP estática); las réplicas lo siguen por su nombre vía el DNS interno de la
VPC.

\begin{lstlisting}[language=none,caption={Creacion del cluster: 1 lider e2-standard-4 + 2 replicas e2-standard-2 (crear-infra.sh, recortado)}]
# Leader: empty tes-leader-host marks it the leader; it gets the static IP.
create_node "${LEADER}" "e2-standard-4" "" "${REPLICA_A}:${HTTP_PORT},${REPLICA_B}:${HTTP_PORT}" \
    --address="${LEADER_IP}"
# Replicas: follow the leader by name; e2-standard-2 each.
create_node "${REPLICA_A}" "e2-standard-2" "${LEADER}" "${LEADER}:${HTTP_PORT},${REPLICA_B}:${HTTP_PORT}"
create_node "${REPLICA_B}" "e2-standard-2" "${LEADER}" "${LEADER}:${HTTP_PORT},${REPLICA_A}:${HTTP_PORT}"
\end{lstlisting}

En cada arranque, \codigo{startup-script.sh} corre como \emph{root} y hace cuatro
cosas: instala el \emph{runtime} \codigo{openjdk-17-jre-headless}; descarga desde
Cloud Storage el \emph{fat jar}, el \emph{dataset} \codigo{alumnos.csv}, la unidad
\codigo{node.service} y la llave \codigo{gcs-key.json}; escribe el archivo de
entorno \codigo{/etc/tesoreria/node.env} con las variables \codigo{TES\_*}
leídas de los metadatos (entre ellas \codigo{TES\_NODE\_ID},
\codigo{TES\_PEERS}, \codigo{TES\_LEADER\_HOST}, \codigo{TES\_JWT\_SECRET} y
\codigo{TES\_BUCKET}); e instala y arranca la unidad \emph{systemd}. Que el mismo
script sirva para todos los roles es lo que mantiene la operación simple: la
diferencia líder/réplica se infiere de \codigo{TES\_LEADER\_HOST}.

La unidad \codigo{deploy/node.service} lanza el \emph{jar} con el puerto HTTP
como primer argumento (\codigo{ExecStart=/usr/bin/java -jar
.../tesoreria-distribuida.jar 8080}), lee el entorno de
\codigo{/etc/tesoreria/node.env} y reinicia \codigo{on-failure} para que un nodo
caído reingrese al clúster automáticamente. Lo \textbf{más importante} de esta
unidad es su política de apagado: usa \codigo{KillSignal=SIGTERM} con
\codigo{TimeoutStopSec=120} para que la JVM ejecute sus \emph{shutdown hooks}. Un
apagado \emph{graceful} con SIGTERM dispara, en una réplica, la escritura del
\emph{checkpoint} exacto (de ahí depende la garantía de reanudar desde la
secuencia precisa) y, en el líder, el drenado de la cola del \emph{journal}
durable. Por eso conviene distinguir reiniciar el \emph{proceso} de reiniciar la
\emph{VM}: matar el proceso con SIGTERM (o \codigo{systemctl stop}) preserva el
estado y el \emph{checkpoint}; nunca debe usarse \codigo{kill -9}, que cortaría la
escritura del \emph{checkpoint} a media subida.

\begin{lstlisting}[language=none,caption={Apagado graceful que preserva el checkpoint (node.service, recortado)}]
KillSignal=SIGTERM
TimeoutStopSec=120
\end{lstlisting}

El ciclo de vida del despliegue se cubre con scripts cortos para no gastar
cómputo cuando no se usa: la infraestructura normalmente está apagada y el día de
la revisión basta \codigo{deploy/encender.sh} para arrancar las cuatro VMs (los
tres nodos más el \codigo{generador}) y esperar a que el líder cargue las 820k
cuentas y sirva \codigo{/api/stats}, sondeándolo con reintentos de \codigo{curl}.
Al terminar, \codigo{deploy/apagar.sh} detiene las VMs sin borrarlas (pausa el
gasto) y \codigo{deploy/destruir-infra.sh} elimina todo, incluidos discos e IP
estática, para no dejar ningún costo. La IP estática del líder no cambia al
apagar y encender, así que el punto de entrada único es estable.

\subsection{Operacion el dia del concurso}

El guion de la demo en vivo está en \codigo{deploy/RUNBOOK.md}, organizado como
\textbf{encender $\rightarrow$ demostrar $\rightarrow$ apagar}. El paso 0 levanta
el clúster con \codigo{encender.sh} y abre el \emph{dashboard} en
\codigo{http://34.67.240.245:8080/} (la IP estática del líder). El paso 1 ejercita
a mano los cuatro \emph{endpoints} con JWT (\codigo{/api/register},
\codigo{/api/login}, \codigo{/api/accounts/\{id\}} y
\codigo{/api/transactions/transfer}) para mostrar sus formas exactas, y el paso 2
demuestra la invariante de conservación: \codigo{totalBalance} en
\codigo{/api/stats} debe quedar constante en \codigo{410113948069.86} antes y
después de transferir.

El corazón de la evaluación es la tolerancia a fallos (pasos 3 y 4). Se demuestra
el ``mata-proceso'' apagando primero \codigo{nodo-3} y luego \codigo{nodo-2} con
\codigo{gcloud compute instances stop}: el líder \textbf{sigue sirviendo} aunque
caigan las dos réplicas. Después se revive una réplica y se observa en su
\emph{log} que reanuda desde su secuencia \emph{exacta} (no desde cero), gracias
al \emph{checkpoint} que escribió al apagarse con SIGTERM, y que converge al líder
(mismo \codigo{lastTxId} y mismo \codigo{totalBalance}). El \emph{dashboard}
muestra todo el tiempo, por nodo, su estado, \#cuentas, saldo total,
\#transferencias, última tx y uso de CPU/RAM/disco, conservando los saldos.

\begin{lstlisting}[language=none,caption={Revivir una replica y verificar reanudacion por secuencia (RUNBOOK paso 4, recortado)}]
$GC compute instances start nodo-2 --zone=$ZONE     # revive
# en su log se ve que reanuda desde su secuencia EXACTA (no desde 0):
#   -> [checkpoint] restored checkpoint/nodo-2.json at watermark <N>
#   -> [node] replica resuming at sequence <N>
\end{lstlisting}

La medición del concurso se automatiza con \codigo{deploy/concurso.sh}, que corre
el \codigo{LoadDriver} \textbf{desde la VM generador} (baja latencia contra la IP
interna del líder) en tres escenarios encadenados, apagando una réplica entre uno
y otro: escenario 1 con el clúster completo (3 nodos), escenario 2 con líder $+$ 1
réplica y escenario 3 solo con el líder. Cada escenario verifica la invariante
(\codigo{CONSISTENT}) e imprime lecturas, transferencias y su
\codigo{score = transferencias*4 + lecturas}; el script suma los tres
\emph{scores} en un total. Que las tres corridas terminen \codigo{CONSISTENT} aun
cruzando apagados de réplicas es la prueba operativa de que el clúster no pierde
ni inventa dinero bajo fallos.

\begin{lstlisting}[language=none,caption={Un escenario del concurso sobre el generador (concurso.sh, recortado)}]
out="$(SSH_GEN "java -cp /opt/tesoreria-distribuida.jar mx.ipn.escom.tesoreria.loadtest.LoadDriver ${LEADER_INT} 8080 ${SECONDS_RUN} ${THREADS} 1 820000 ${label} 2>&1")"
echo "${out}" | grep -E 'CONSISTENT|INCONSISTENT|successful|score' >&2
\end{lstlisting}

Para cerrar, el RUNBOOK documenta la recuperación ante desastre total: como el
estado durable vive en Cloud Storage, un arranque en frío de los nodos (incluso
tras destruir las VMs y volver a crearlas) reconstruye el banco desde el bucket.
La limpieza final del día es \codigo{deploy/apagar.sh} para solo pausar el gasto,
o \codigo{deploy/destruir-infra.sh} para borrar discos e IP estática y no dejar
ningún costo. El RUNBOOK advierte además que las réplicas siempre se apaguen con
SIGTERM y nunca con \codigo{kill -9}, para no truncar la escritura del
\emph{checkpoint}.
